\documentclass[11pt]{article}
\usepackage[margin=1in]{geometry}
\usepackage{amsmath,amssymb}
\pagestyle{empty}

\title{Deriving the Chebyshev Distance Using Limits}
\author{}
\date{}

\begin{document}
\maketitle

For a vector $z=(z_1,\ldots,z_n)\in\mathbb{R}^n$, with $n\geq 1$, the
$\ell_p$ norm is
\[
  \|z\|_p = \left(\sum_{i=1}^n |z_i|^p\right)^{1/p},
  \qquad 1\leq p<\infty.
\]
We will show that its limit as $p\to\infty$ is the largest component magnitude.

\paragraph{Simple proof.}
Let $M=\max_{1\leq i\leq n}|z_i|$. At least one term in the sum equals
$M^p$, and every term is at most $M^p$. Hence
\[
  M^p \leq \sum_{i=1}^n |z_i|^p \leq nM^p.
\]
Taking the $p$th root gives
\[
  M \leq \|z\|_p \leq n^{1/p}M.
\]
Since
\[
  \lim_{p\to\infty}n^{1/p}
  =\lim_{p\to\infty}e^{(\log n)/p}=1,
\]
the squeeze theorem yields
\[
  \boxed{\lim_{p\to\infty}\|z\|_p
  =\max_{1\leq i\leq n}|z_i|
  =\|z\|_\infty.}
\]
This argument also covers $z=0$, when both bounds are zero.

\paragraph{Chebyshev distance.}
Applying the result to $z=x-y$ gives the $\ell_\infty$ (Chebyshev) distance:
\[
  \boxed{d_\infty(x,y)
  =\lim_{p\to\infty}\left(\sum_{i=1}^n|x_i-y_i|^p\right)^{1/p}
  =\max_{1\leq i\leq n}|x_i-y_i|.}
\]

\paragraph{Note.}
The $\ell_\infty$ norm selects a component with maximum magnitude and returns
that magnitude, not its sign or index. For distance, it selects the largest
absolute coordinate difference. Ties do not affect the value.

\end{document}
